%!TEX root = tfg.tex
\chapter{Ejecución de los Planes de Gestión}

\begin{quotation}[Founding Father]{John Adams (1735--1826)}
Facts are stubborn things.
\end{quotation}

\begin{abstract}
Este capítulo describe cómo los planes de gestión definidos previamente se trasladan a la
ejecución del proyecto. Su finalidad no es redefinir la planificación, sino explicar cómo se
aplican sus mecanismos de control durante el desarrollo de \textit{Via Inferni}.
\end{abstract}

\section{Ejecución del plan de gestión del alcance}
La ejecución del alcance se articula mediante FDD \cite{palmer2002fdd}. El modelo general
define los grandes dominios del videojuego y la lista de funcionalidades concreta el trabajo
que puede planificarse, implementarse y validarse en iteraciones.

\subsection{Modelo general}
El modelo general proporciona una visión conceptual del producto. Aunque FDD lo define
formalmente como una aproximación inicial a la arquitectura del sistema, en este proyecto se
utiliza con un alcance más funcional. No se pretende fijar una arquitectura antes de las
iteraciones, sino agrupar las grandes áreas de trabajo y las funcionalidades principales que
después serán diseñadas e implementadas de forma incremental. En \textit{Via Inferni}, el
sistema se organiza en cuatro dominios:

\begin{itemize}
    \item \textbf{Generación y estructura del mundo:} construcción de niveles, distribución
    de salas, conectividad y progresión espacial.
    \item \textbf{Jugador, enemigos y combate:} entidad jugable, enemigos, desplazamiento,
    acciones ofensivas, recepción de daño y comportamiento de oposición.
    \item \textbf{Interfaz y experiencia de uso:} cámara, minimapa, \textit{Head-Up Display}
    y menús que comunican el estado de la partida.
    \item \textbf{Progresión y gestión de recursos:} objetos, inventario y elementos que
    aportan persistencia, toma de decisiones y variedad a cada partida.
\end{itemize}

La relación entre estos dominios es dependiente: la generación crea el espacio jugable; el
jugador y los enemigos actúan sobre él; la interfaz permite interpretar lo que ocurre; y la
progresión aporta continuidad a la experiencia.

\subsection{Lista de funcionalidades}
A partir del modelo general se construye una lista de funcionalidades agrupadas por dominio.
Esta lista permite seleccionar trabajo para cada iteración y mantener trazabilidad entre
alcance, implementación y validación.

\subsubsection{Generación y estructura del mundo}
\begin{itemize}
    \item \textbf{Generación procedural:} creación automática de niveles, salas y
    conectividad.
    \item \textbf{Lógica de círculos:} progresión por los niveles temáticos del Infierno.
    \item \textbf{Sala tienda:} punto de compra y preparación dentro de la \textit{run}.
    \item \textbf{Sala secreta y tesoro:} contenido oculto y recompensas de exploración.
\end{itemize}

\subsubsection{Jugador, enemigos y combate}
\begin{itemize}
    \item \textbf{Jugador:} entidad controlable e integrada en el nivel generado.
    \item \textbf{Cámara del jugador:} seguimiento visual para facilitar navegación.
    \item \textbf{Dualidad de personajes:} alternancia y persistencia compartida entre
    Dante y Virgilio.
    \item \textbf{Enemigos:} incorporación de oposición básica en las salas.
    \item \textbf{Sistema de combate:} reglas para atacar, recibir daño y resolver
    enfrentamientos.
    \item \textbf{Sistema de enemigos:} comportamiento, estados y escalado de dificultad.
\end{itemize}

\subsubsection{Interfaz y experiencia de uso}
\begin{itemize}
    \item \textbf{Minimapa:} representación simplificada del nivel.
    \item \textbf{Head-Up Display:} información esencial de vida, recursos y progreso.
    \item \textbf{Menú de pausa:} interrupción controlada de la partida.
    \item \textbf{Menú principal:} acceso inicial a partida y opciones principales.
    \item \textbf{Actualización de sprites:} mejora visual y coherencia estética.
    \item \textbf{Victoria y derrota:} flujos de cierre de partida.
\end{itemize}

\subsubsection{Progresión y recursos}
\begin{itemize}
    \item \textbf{Gestión del inventario:} almacenamiento y administración de recursos.
    \item \textbf{Sistema de objetos:} elementos interactivos, pasivos o consumibles que
    modifican la partida.
\end{itemize}

\section{Ejecución del plan de gestión del cronograma}
El cronograma se ejecuta mediante tres fases y una fase central compuesta por iteraciones.
Los hitos se utilizan como puntos de control para revisar desviaciones, reajustar carga de
trabajo y decidir si una funcionalidad debe cerrarse, dividirse o posponerse.

\subsection{Aplicación de las fases}
La Fase 1 prepara el marco de desarrollo: objetivos, repositorio, backlog inicial,
estructura documental y herramientas. La Fase 2 concentra el desarrollo técnico mediante
iteraciones cerradas. La Fase 3 queda reservada para consolidar el producto, revisar la
memoria y preparar la entrega final.

\subsection{Replanificación iterativa}
Durante la ejecución, el cronograma no se trata como una secuencia rígida. Si una
funcionalidad resulta más compleja de lo previsto, se divide en entregas menores o se
desplaza parcialmente a una iteración posterior. Este criterio permite mantener incrementos
verificables sin bloquear un ciclo completo por una única tarea.

\subsection{Uso de hitos y margen temporal}
Al cierre de cada iteración se revisa qué se ha completado, qué queda pendiente y qué debe
ajustarse. El margen temporal previsto en la planificación permite absorber desviaciones,
resolver incidencias y preparar el cierre documental sin comprometer la entrega final.

\section{Ejecución del plan de gestión de comunicaciones}
La comunicación diaria se apoya en Discord, mientras que GitHub concentra la trazabilidad
formal del trabajo \cite{githubIssuesProjects}. El tablero Kanban permite distinguir tareas
pendientes, en curso, en revisión y cerradas.

\subsection{Comunicación operativa}
Discord se utiliza para resolver bloqueos rápidos, coordinar decisiones de bajo impacto y
mantener continuidad entre sesiones de trabajo. Este canal evita convocar reuniones formales
para cada ajuste operativo.

\subsection{Trazabilidad formal}
GitHub registra incidencias, commits, ramas, revisiones y evolución del backlog. Esta
información permite reconstruir la secuencia real de trabajo y alinear lo implementado con
lo descrito en la memoria.

\subsection{Reuniones de iteración}
Las reuniones de planificación seleccionan funcionalidades, revisan dependencias y fijan
criterios de aceptación. Las reuniones de cierre validan resultados, documentan incidencias,
actualizan riesgos y preparan la siguiente iteración.

\section{Ejecución del plan de gestión de riesgos}
El plan de riesgos se ejecuta como un instrumento vivo. Los riesgos identificados en la
planificación se revisan durante las iteraciones y se actualizan si aparecen nuevas
incertidumbres.

\subsection{Riesgos técnicos}
La generación procedural, la integración progresiva de sistemas y la dependencia de Unity se
vigilan especialmente. Su control se apoya en prototipos tempranos, integración frecuente,
validación en contexto real de juego y estabilidad de versiones.

\subsection{Riesgos de plazo}
El sobre-alcance, las estimaciones optimistas y los bloqueos por dependencias se controlan
mediante priorización, división de funcionalidades, uso del backlog y reserva de margen
temporal. Las desviaciones relevantes se documentan al cierre de la iteración en la que se
detectan.

\section{Ejecución del plan de gestión de la calidad}
La calidad se ejecuta mediante la aplicación repetida de la \textit{Definition of Done}. Una
funcionalidad no se considera cerrada solo por existir en código; debe compilar, integrarse,
no introducir regresiones, revisarse y validarse mediante pruebas jugables.

\subsection{Control por iteración}
Cada iteración incorpora validaciones técnicas y \textit{playtests} definidos mediante pasos
reproducibles. Cuando una prueba detecta una incidencia, esta se registra, se corrige si es
crítica para la iteración o se incorpora al backlog si puede abordarse posteriormente sin
comprometer el incremento.

\subsection{Calidad acumulativa}
El control de calidad tiene carácter acumulativo. Cada nueva funcionalidad debe funcionar
por sí misma y mantener estable lo ya implementado. De este modo, el proyecto evita acumular
deuda técnica no visible hasta la fase de cierre.

\section{Ejecución del plan de gestión de costes}
La ejecución del plan de costes se refleja mediante una estimación consolidada del valor
económico del proyecto. El mayor peso corresponde al tiempo de trabajo de los dos
desarrolladores, complementado con la amortización del hardware y los costes indirectos.

\subsection{Costes de personal}
Se toma como referencia un perfil junior de ingeniería de software en España con un salario
bruto anual de 35.000 \euro, dentro de los rangos publicados para perfiles tecnológicos en
la guía salarial de Manfred de 2026 \cite{manfredGuiaSalarial2026}. Con una dedicación
estimada de 300 horas por integrante y una jornada anual de 1.800 horas, la imputación
equivale a una sexta parte de la anualidad.

\begin{itemize}
    \item 35.000 \euro\ / 6 = 5.833,33 \euro\ por persona.
    \item 5.833,33 \euro\ $\times$ 2 = 11.666,66 \euro.
    \item 11.666,66 \euro\ $\times$ 0,30 = 3.500,00 \euro\ en costes sociales.
\end{itemize}

El coste total de personal, incluyendo costes sociales, asciende a 15.166,66 \euro.

\subsection{Costes materiales e indirectos}
Se considera el uso de dos equipos informáticos valorados en 1.500 \euro\ cada uno. Con una
amortización anual del 25\% y una imputación equivalente a una sexta parte del año, el coste
total de amortización asciende a 125,00 \euro. El coste de software y servicios se considera
nulo, al utilizar herramientas gratuitas, educativas o ya disponibles.

Los costes indirectos se calculan como el 10\% de los costes directos. La suma de costes
directos es de 15.291,66 \euro, por lo que los costes indirectos ascienden a
1.529,17 \euro.

\subsection{Resultado global}
El valor económico total estimado del proyecto asciende a 16.820,83 \euro. Aunque no
representa un desembolso real, permite dimensionar el esfuerzo realizado con criterios
profesionales. La Tabla~\ref{tab:resumen_presupuesto} recoge el desglose del presupuesto
consolidado del proyecto.

\begin{table}[!htbp]
    \centering
    \begin{coolTable}{lr}{2}
    {Resumen del presupuesto del proyecto}
    \textbf{Costes de personal (salarios brutos)} & 11.666,66 \euro \\
    \textbf{Costes sociales (30\% sobre bruto)} & 3.500,00 \euro \\
    \textbf{Costes materiales (amortización hardware)} & 125,00 \euro \\
    \textbf{Costes materiales (software y servicios)} & 0,00 \euro \\
    \midrule
    \textbf{Total costes directos} & 15.291,66 \euro \\
    \textbf{Costes indirectos (10\%)} & 1.529,17 \euro \\
    \midrule
    \textbf{TOTAL DEL PROYECTO} & \textbf{16.820,83 \euro} \\
    \end{coolTable}
    \caption{Resumen detallado de la inversión económica total.}
    \label{tab:resumen_presupuesto}
\end{table}
